% §5.4 RQ4 alert shifts and telemetry ablation, Observe lane (O-12, #48). Input by 05-results.tex
% through \lanesectionor{observe-rq4}{05-rq4-shift-ablation}; it replaces the lead's draft and keeps
% its tables, figures and \resnum keys. Sources: docs/04 §2 (H4a, H4b), §5.2, §6; engines 01 §3.7,
% §3.10, 02 §2 (window independence), 03 §3.7; bench/results/not_run.yaml at 25 Nov 2026 (#532).
% No result number is typed here; every value is a \resnum key or a generated table.
\paragraph{Alert-time shifts}
A shift moves the alert time and nothing else.
Core's alert step adds it to the oracle time and recomputes the two windows, and the engines cannot tell a shifted alert from an unshifted one.
The series, the log templates and their identifiers do not depend on the alert (\S\ref{sec:pipeline}), so a shift changes only what is computed per window: which points count as baseline, the edge and critical-path tables, the per-window template counts, and through them the anomaly scores and onsets.
Two parts of the design bound its effect.
An alert that is 120\,s early still has the fault inside the incident window, which runs to eight minutes after the alert.
An alert that is 120\,s late puts the first two minutes of the fault at the very start of the incident window, and the incident onset is clamped to no earlier than 60\,s before the alert, so onsets found earlier are pulled forward.
We therefore expected late alerts to cost more than early ones; that expectation is not a pre-registered hypothesis.

\resultfigure*{fig:shift}{AC@1 per RE2 system against the shift of the alert time ($-120$ to $+120$\,s; $\pm 30$\,s exploratory) for N-Sigma, BARO, alert counting, the fused slate and \code{veyltrace:off}; \code{veyltrace:brief} has a 0\,s point only, because its shifted runs were cut.}{re2-oracle-rq1-v1 and re2-oracle-rq2-v1 (0 s), re2-shift-{m120,m60,m30,p30,p60,p120}-v1}{fig7-shift}
\resulttable*{tab:shift}{Alert-shift robustness (H4a): AC@1, AC@3 and Retrieval@K per shift, system and method.}{re2-oracle-rq1-v1, re2-oracle-rq2-v1, re2-shift-{m120,m60,m30,p30,p60,p120}-v1}{t10-rq4-shift}
H4a holds if AC@1 drops by at most 0.10 per RE2 system under $\pm 60$ and $\pm 120$\,s shifts.
Over those four shifts, the largest drop of \code{veyltrace:off} is \resnum{re2-oracle-rq1-v1}{ob/veyltrace:off/ac_at_1/max_drop/vs:re2-shift-m120-v1,re2-shift-m60-v1,re2-shift-p60-v1,re2-shift-p120-v1} on OB, \resnum{re2-oracle-rq1-v1}{ss/veyltrace:off/ac_at_1/max_drop/vs:re2-shift-m120-v1,re2-shift-m60-v1,re2-shift-p60-v1,re2-shift-p120-v1} on SS and \resnum{re2-oracle-rq1-v1}{tt/veyltrace:off/ac_at_1/max_drop/vs:re2-shift-m120-v1,re2-shift-m60-v1,re2-shift-p60-v1,re2-shift-p120-v1} on TT (Figure~\ref{fig:shift}, Table~\ref{tab:shift}).
The six shifted runs of the AI arm were pre-registered and are reported as not run (Appendix~\ref{sec:appendix}), so H4a is tested for the deterministic path only.
\leadq{the lead's draft of this caption and table listed \code{re2-shift-*-brief-v1}; all six are in \code{not_run.yaml} since 25 Nov (\#532), so I removed them here. Check that R-17's \code{fig7-shift} and \code{t10-rq4-shift} agree.}

\paragraph{Missing telemetry}
A dropped source is removed at ingest: its table is written with no rows and its flag is cleared, exactly as for a source that the dataset never had, so no engine can treat the two differently (\S\ref{sec:pipeline}).
Downstream, a ranker whose channel is missing abstains from the fused score (\S\ref{sec:slate}), and a trial that needs the channel is untestable, which is not a pass (\S\ref{sec:verifier}); the design intent is that missing telemetry lowers the status of a verdict before it lowers its correctness, which is what H4b measures.
Two cells are controls by construction.
RCAEval has no change events, so dropping events removes nothing, and that cell should reproduce the reference cell exactly; a difference there would be a defect in the hook, not an effect of telemetry.
RE2-SS has no traces, so the same holds for dropping traces on SS.

\resultfigure*{fig:ablation}{Telemetry ablation on RE2: AC@1 of \code{veyltrace:off} per cell (no source dropped; drop metrics, logs, traces or events; metrics only), one panel per system.}{re2-oracle-rq1-v1 (no source dropped), re2-drop-metrics-v1, re2-drop-logs-v1, re2-drop-traces-v1, re2-drop-events-v1, re2-metrics-only-v1}{fig8-ablation}
\resulttable{tab:ablation}{Telemetry ablation (H4b): accuracy and the share of correct \st{CONFIRMED} verdicts per cell, system and method.}{re2-oracle-rq1-v1, re2-drop-metrics-v1, re2-drop-logs-v1, re2-drop-traces-v1, re2-drop-events-v1, re2-metrics-only-v1}{t11-rq4-ablation}
H4b holds if, in every ablation cell, at least 90\% of \st{CONFIRMED} verdicts are correct.
Table~\ref{tab:ablation} gives that share per cell and system, and Figure~\ref{fig:ablation} the accuracy.
\leadq{kept from the lead's draft: \code{bench_summary} 1.1 has no \code{selective_accuracy} entry, so H4b needs a contract change before its numbers can be generated. The mix of statuses per cell, which docs/11 §5 plans for this figure, needs a summary entry too.}
